% Part: normal-modal-logic
% Chapter: axioms-systems
% Section: normal-logics

\documentclass[../../../include/open-logic-section]{subfiles}

\begin{document}

\olfileid{nml}{prf}{nor}

\olsection{Logika Modal Normal}

Ora saben himpunan !!{formula}s modal gampang dicirèkaké
minangka himpunan !!{formula}s kang bisa diderivasi saka sawijining
himpunan aksioma. Kita pengin logika modal nduweni sipat kang trep.
Kaping pisan, kabeh kang bisa kita !!{derive} ing logika proposisional
klasik mesthi kudu tetep !!{derivable}, kanthi nggatekake yen
!!{formula}s saiki uga bisa ngemot
\iftag{prvBox}{$\Box$\iftag{prvDiamond}{
    lan~}{}}{}\iftag{prvDiamond}{$\Diamond$}{}. Kanggo iku,
kita mbutuhake logika modal ngemot kabeh instansi tautologis lan
ditutup marang modus ponens.

\begin{defn}
  Sawijining \emph{logika modal} yaiku himpunan~$\Sigma$ saka
  !!{formula}s modal kang
  \begin{enumerate}
  \item ngemot kabeh tautologi, lan
  \item ditutup marang substitusi, yaiku yen $!A \in \Sigma$, lan
    $!D_1$, \dots, $!D_n$ iku !!{formula}s, mula
    \[
    \SSubst{!A}{\subst{!D_1}{p_1}, \dots, \subst{!D_n}{p_n}} \in \Sigma,
    \]
    \item ditutup marang \emph{modus ponens}, yaiku yen $!A$ lan $!A
      \lif !B \in \Sigma$, mula $!B \in \Sigma$.
  \end{enumerate}
\end{defn}

Supaya bisa nggunakake semantik relasional kanggo logika
modal, kita uga kudu mbutuhake yen kabeh !!{formula}s kang valid ing
kabeh model modal kalebu. Nyatane syarat iki wis kacukupan yen kabeh
instansi \Ax{K}\iftag{prvDiamond}{ lan~\Dual{}}{} iku
!!{derivable}, lan saben !!a{formula}~$!A$ kang !!{derivable} uga
ndadekake~$\Box !A$ bisa diderivasi. Logika modal kang nyukupi
syarat-syarat iki diarani \emph{normal}. (Mesthi ana uga logika modal
kang ora normal, nanging model relasional lumrah ora cocog kanggo
logika kasebut.)

\begin{defn}
  Logika modal $\Sigma$ iku \emph{normal} yen ngemot
  \begin{align*}
    \tag{\Ax{K}} & \Box(p \lif q) \lif (\Box p \lif \Box q),
    \iftag{prvDiamond}{\\
      \tag{\Dual} & \Diamond p \liff \lnot\Box\lnot p}{}
  \end{align*}
  lan ditutup marang \emph{nesesitasi}, yaiku yen $!A \in
  \Sigma$, mula $\Box !A \in \Sigma$.
\end{defn}

Gatekna yen implikasi tautologis cukup ``rinci'' kanggo
njaga \emph{kabeneran ing sawijining donya}, nanging aturan \Nec{}
mung njaga \emph{kabeneran ing sawijining model} (lan mula uga
validitas ing sawijining bingkai utawa kelas bingkai).

\begin{prop}\ollabel{prop:rk}
  Saben logika modal normal ditutup marang aturan \RK,
  \begin{prooftree}
    \AxiomC{$!A_1 \lif (!A_2 \lif \cdots (!A_{n-1} \lif !A_n)\cdots)$}
    \RightLabel{\RK}
    \UnaryInfC{$\Box!A_1 \lif (\Box!A_2 \lif \cdots (\Box!A_{n-1}
      \lif \Box!A_n)\cdots).$}
  \end{prooftree}
\end{prop}

\begin{proof}
  Kanthi induksi marang~$n$: Yen $n = 1$, aturan iku mung \Nec,
  lan saben logika modal normal ditutup marang \Nec.

Saiki upamane asil iki bener kanggo $n-1$; kita tuduhake
yen uga bener kanggo~$n$.

Upamane
  \begin{align*}
  & !A_1 \lif (!A_2 \lif \cdots (!A_{n-1} \lif !A_n)\cdots) \in \Sigma
  \intertext{Miturut hipotesis induksi, kita duwe}
  & \Box!A_1 \lif (\Box!A_2 \lif \cdots \Box(!A_{n-1} \lif !A_n)\cdots)
  \in \Sigma
  \intertext{Amarga $\Sigma$ logika modal normal, kabeh instansi
    saka~$\Ax{K}$ kalebu, mligine}
  & \Box(!A_{n-1} \lif !A_n) \lif (\Box!A_{n-1} \lif \Box!A_n) \in \Sigma
  \intertext{Kanthi modus ponens lan instansi tautologis kang cocog, kita entuk}
  & \Box!A_1 \lif (\Box!A_2 \lif \cdots (\Box!A_{n-1}
  \lif \Box!A_n)\cdots) \in \Sigma.
  \end{align*}
\end{proof}

\begin{prop}\ollabel{prop:notDiamondBot}
  Saben logika modal normal $\Sigma$ ngemot~$\lnot\Diamond\lfalse$.
\end{prop}

\begin{prob}
  Buktekna \olref[nml][prf][nor]{prop:notDiamondBot}.
\end{prob}

\begin{prop}
  Upamane $!A_1$, \dots, $!A_n$ iku !!{formula}s. Mula ana logika
  modal normal paling cilik $\Sigma$ kang ngemot kabeh instansi
  saka $!A_1$, \dots, $!A_n$.
\end{prop}

\begin{proof}
  Kanggo $!A_1$, \dots, $!A_n$, temtokna $\Sigma$ minangka irisan
  kabeh logika modal normal kang ngemot kabeh instansi saka
  $!A_1$, \dots, $!A_n$. Kulawarga kang diiris ora kosong, amarga
  $\Frm[L]$, himpunan kabeh !!{formula}s, uga logika modal normal
  kaya mangkono.
\end{proof}

\begin{defn}
Logika modal normal paling cilik kang ngemot $!A_1$, \dots, $!A_n$
diarani \emph{sistem modal} lan dilambangake $\Log{K} !A_1 \dots
!A_n$. Logika modal normal paling cilik dilambangake~\Log{K}.
\end{defn}

\end{document}
